% Tabellen-Einstellungen für dieses Formular (tabularx-Tabellen werden durch
% ltablex zu longtables und können über Seiten umbrechen)
\keepXColumns
\renewcommand{\theHtable}{ki.\arabic{table}} % eindeutige Hyperref-Anker (sonst Konflikt mit Tabelle A.1)
\setlength{\extrarowheight}{2pt}
\newcolumntype{L}{>{\raggedright\arraybackslash}X}

% --------------------------------------------------
% Erklärung zur Nutzung von generativen KI-Tools
% --------------------------------------------------

\section*{Erklärung zur Nutzung von generativen KI-Tools in Hochschulprüfungen}

\noindent
(In Anlehnung an \citealp{gimpel2023kognition})\\[0.5\baselineskip]
Die Erklärung soll für Seminar-, Projekt-, Bachelor- und Masterarbeiten verwendet werden, um die Nutzung von generativen KI-Tools und die eigene wissenschaftliche Leistung zu erläutern.

\vspace{1em}

\noindent
\textbf{Titel der Arbeit} \rule{0.7\textwidth}{0.4pt}\\[0.6\baselineskip]
\textbf{Abgabedatum} \rule{0.7\textwidth}{0.4pt}

\vspace{2em}

\subsection*{1. Verantwortungsvoller Umgang mit KI-Tools}

Bestätigen Sie die folgenden Aussagen.
\begingroup
\small
\begin{tabularx}{\textwidth}{|L|L|}
\hline
\textbf{Aussagen zum verantwortungsvollen Umgang mit KI-Tools} &
\textbf{Bestätigung: Ja} \\
\hline
Ich bin/wir sind über die Möglichkeiten und Grenzen der von mir/uns
verwendeten generativen KI-Tools informiert. &
 \\
\hline
Ich/wir habe(n) mich/uns vergewissert, dass die von den KI-Tools
gelieferten Ergebnisse korrekt sind, oder dass ich/wir sie korrigiert
habe(n). &
 \\
\hline
Ich/wir erkenne(n) an, dass die Verantwortung für die Arbeit bei
den Verfassenden liegt, nicht bei den KI-Tools oder einer anderen
Person. &
 \\
\hline
\end{tabularx}
\endgroup


\subsection*{2. Detaillierte Tätigkeiten, für die KI-Tools genutzt wurden}
Geben Sie an, welches KI-Tool Sie in welchem Umfang für welche Tätigkeit genutzt haben.
\begingroup
\small

\begin{tabularx}{\textwidth}{|L|L|L|L|}
\hline
\textbf{Aktivität} &
\textbf{Beschreibung} &
\textbf{Verwendete KI-Tools}\textsuperscript{*} \newline \scriptsize (falls zutreffend) &
\textbf{Beschreibung der Nutzung der KI-Tools} \newline \scriptsize (Art der Nutzung, betroffene Abschnitte der Arbeit, etc.) \\
\hline
\endhead
\hline
\endfoot
\textbf{Ideenfindung \& Konzeptualisierung} &
– Entwicklung von Ideen, Forschungszielen, Zielsetzungen und Forschungsfragen \newline
– Identifikation und Definition relevanter Konzepte &
&
\\
\hline
\textbf{Literaturrecherche und -analyse} &
– Suche nach relevanter Literatur \newline
– Überprüfung von potenziell relevanter Literatur \newline
– Zusammenfassungen von relevanter Literatur &
&
\\
\hline
\textbf{Methodik} &
– Suche nach einer geeigneten Methodik \newline
– Entwicklung und Anpassung der Methodik an die Forschungsfrage(n) &
&
\\
\hline
\textbf{Coding} &
– Erstellen und Dokumentieren von Code, Algorithmen, Software \newline
– Testen und Debuggen von bestehendem Code, Algorithmen, Software \newline
– Verständnis von bestehendem Code, Algorithmen, Software &
&
\\
\hline
\textbf{Datenerhebung und -analyse} &
– Erhebung von Primär- oder Sekundärdaten \newline
– Qualitative Datenanalyse (einschließlich Zusammenfassungen und Coding) \newline
– Quantitative Datenanalyse (einschließlich statistische Auswertung) \newline
– Mathematische, computergestützte oder andere formale Techniken zur Modellierung, Simulation und Analyse &
&
\\
\hline
\textbf{Interpretation und Validierung} &
– Interpretation der Ergebnisse \newline
– Ableitung von Implikationen für Forschung und Praxis \newline
– Überprüfung der allgemeinen Replikation/Reproduzierbarkeit der Ergebnisse und anderer Forschungsergebnisse &
&
\\
\hline
\textbf{Strukturierung und Planung des Textes} &
– Gliederung der Arbeit \newline
– Gliederung von Abschnitten der Arbeit (z.\,B. Stichpunktliste pro Abschnitt) &
&
\\
\hline

\textbf{Generierung des Textes} &
– Generierung von Text zu verschiedenen Themen in verschiedenen Abschnitten der Arbeit (einschließlich Titel und Zusammenfassungen) &
&
\\
\hline
\textbf{Übersetzen von Text} &
– Übersetzung selbst verfasster Texte \newline
– Übersetzung von anderen verfasster Texte &
&
\\
\hline
\textbf{Überprüfung \& Editierung des Textes} &
– Kritische Überprüfung, Feedback oder Überarbeitung des Inhalts, der Organisation oder der Grammatik der Arbeit \newline
– Korrekturlesen \newline
– Umformulierung oder Paraphrasierung von Text \newline
– Kürzung/Erweiterung des Textes &
&
\\
\hline
\textbf{Präsentation} &
– Strukturierung einer Präsentation über die Arbeit \newline
– Füllen einer Präsentation über die Arbeit mit Inhalt &
&
\\
\hline
\textbf{Quellenverwaltung} &
– Erstellen der Referenzliste \newline
– Formatierung der Referenzen &
&
\\
\hline
\textbf{Weitere Aktivitäten} & & & \\
\end{tabularx}
\endgroup
\vspace{2em}
\noindent\scriptsize
\textsuperscript{*}\,In Spalte 3 (Verwendete KI-Tools) nennen Sie bitte die konkrete Bezeichnung und Version der eingesetzten Tools (z.\,B. ChatGPT (Version xx), Perplexity, DeepL Write usw.). In Spalte 4 (Beschreibung der Nutzung) beschreiben Sie knapp, wie das Tool eingesetzt wurde (Art der Nutzung) und auf welche Teile der Arbeit sich die Nutzung bezog (z.\,B. „Überarbeitung der Diskussion“, „Entwurf der Gliederung“, „erste Textfassung des Abschnitts 3.2“).
\vspace{2em}

\subsection*{3. Ergänzende Angaben}

Gibt es sonst noch etwas, das Sie bezüglich der Nutzung von KI-Tools angeben wollen, welches Ihre eigenständige, definierbare Leistung bei der Erstellung dieser Arbeit hervorheben oder einschränken würde?

\vspace{1\baselineskip}
\noindent\rule{\textwidth}{0.4pt}\\[0.8\baselineskip]
\noindent\rule{\textwidth}{0.4pt}\\[0.8\baselineskip]
\noindent\rule{\textwidth}{0.4pt}

\vspace{2em}

\subsection*{4. Unterschrift(en) der Verfassenden}

\begingroup
\small

\begin{tabularx}{\textwidth}{|X|X|X|X|}
\hline
\textbf{Name} &
\textbf{Immatrikulations Nr.} &
\textbf{Datum} &
\textbf{Unterschrift} \\
\hline
 & & & \\
\hline
 & & & \\
\hline
 & & & \\
\hline
 & & & \\
\hline
 & & & \\
\hline
\end{tabularx}

\endgroup

\vspace{2em}
